% =========================================================
% Completeness of the Real Numbers — Cross-Reference
%
% The full development (Completeness Axiom, Nested Interval Property,
% Archimedean Property, Floor Lemma, Density of Q, Existence of sqrt)
% is given in the Bounding chapter:
%
%   volume-iii/analysis/bounding/notes/notes-completeness.tex
%
% That chapter precedes Convergence in the curriculum, and the
% theorems there are proved in full with stubs in
%   volume-iii/analysis/bounding/proofs/notes/
%
% This file replaces the prior duplicate of that content with a
% forward reference.
% =========================================================

\subsection{Completeness --- Reference to Bounding Chapter}

\begin{tcolorbox}[colback=gray!6, colframe=gray!40, arc=2pt,
  left=6pt, right=6pt, top=4pt, bottom=4pt,
  title={\small\textbf{These results are proved in the Bounding chapter}},
  fonttitle=\small\bfseries]
The Completeness Axiom and its consequences --- Nested Interval Property,
Archimedean Property, Floor Lemma, Density of $\mathbb{Q}$, Density of the
Irrationals, and Existence of Square Roots --- are developed fully in the
\textbf{Bounding} chapter (\texttt{volume-iii/analysis/bounding/}).

\medskip

For proofs, motivation, and diagrams consult the Bounding chapter notes.
The proof stubs are in
\texttt{volume-iii/analysis/bounding/proofs/notes/}.

\medskip
\textbf{What you should already know entering this chapter:}
\begin{itemize}
  \item The LUB property: every nonempty bounded-above $S \subseteq \mathbb{R}$
        has $\sup S \in \mathbb{R}$.
  \item The $\varepsilon$-characterization: $s = \sup A$ iff $s$ is an upper
        bound and for every $\varepsilon > 0$ there exists $a \in A$ with
        $s - \varepsilon < a$.
  \item The Archimedean property: $\forall x \in \mathbb{R}$,
        $\exists n \in \mathbb{N}$ with $n > x$.
  \item Density of $\mathbb{Q}$: between any two reals lies a rational.
\end{itemize}
\end{tcolorbox}

\begin{remark*}[Why completeness matters \emph{here}]
This chapter is where completeness first does visible work.
The Monotone Convergence Theorem asserts that a bounded monotone sequence
has a limit --- but that limit is $\sup\{x_n\}$, which exists only because
of the LUB property.
Without completeness (e.g.\ over $\mathbb{Q}$) the supremum might fail to
exist as an element of the ambient field, and the theorem would fail.
\end{remark*}
